% ==============================================================================
% ICML 2026 MechInterp Workshop Paper (4 pages + references)
% Self-contained file — does not use \input
% ==============================================================================

\documentclass[accepted]{article}

% ---------- Packages ----------
\usepackage[utf8]{inputenc}
\usepackage[T1]{fontenc}
\usepackage{amsmath,amssymb,amsfonts}
\usepackage{graphicx}
\usepackage{booktabs}
\usepackage{enumitem}
\usepackage{hyperref}
\usepackage{url}
\usepackage{xcolor}
\usepackage{natbib}
\usepackage{microtype}
\usepackage[margin=1in]{geometry}

% NOTE: Replace \documentclass and geometry with the official ICML workshop
% style file when available. For submission, switch to:
%   \documentclass{icml2026}
% and remove the geometry package.

% ---------- Custom commands ----------
\newcommand{\sone}{S1-like}
\newcommand{\stwo}{S2-like}
\newcommand{\pzero}{\texttt{P0}}
\newcommand{\pending}{\textcolor{gray}{\textit{[pending, running]}}}

% ---------- Title ----------
\title{What Changes When LLMs Learn to Reason?\\
  Mechanistic Signatures of Cognitive Bias Processing}

\author{%
  \textbf{Bright Liu}\textsuperscript{1}\thanks{Correspondence: \texttt{brightliu@college.harvard.edu}} \quad
  \textbf{[Co-authors]}\textsuperscript{1} \quad
  \textbf{[Advisor]}\textsuperscript{1,2} \\[6pt]
  \textsuperscript{1}Harvard Undergraduate AI Safety (HUSAI) \quad
  \textsuperscript{2}[Department] \\
}

\date{}

\begin{document}

\maketitle

% ============================================================
% Abstract
% ============================================================

\begin{abstract}
Reasoning-trained large language models resist cognitive bias lures that
trip up their standard counterparts, but does this behavioral shift
reflect a genuine change in internal computation?  We present a
contrastive benchmark of 330 items (165 matched conflict/control pairs)
across 7 cognitive bias categories and apply linear probes to residual
stream activations of models that share \emph{identical} architectures
but differ in training.  Behaviorally, Llama-3.1-8B-Instruct falls for
heuristic lures at a 27.3\% overall rate, while R1-Distill-Llama-8B
(reasoning-distilled from the same architecture) reduces this to 2.4\%.
Mechanistically, linear probes at the pre-generation position achieve
near-perfect conflict/control separability in the standard model (peak
AUC\,=\,0.999, layer 14) but consistently lower separability in the
reasoning model (peak AUC\,=\,0.929)---the internal \sone{}/\stwo{}
boundary is \emph{blurred}, not sharpened, by reasoning training.  A
within-model toggle using Qwen~3-8B with and without thinking mode
provides a further control \pending{}.  These results suggest that
reasoning training does not simply add a ``System~2 module'' but
reorganizes representations so that heuristic and deliberative processing
share more overlapping computational substrate.
\end{abstract}

% ============================================================
% 1. Introduction
% ============================================================

\section{Introduction}
\label{sec:intro}

Kahneman's dual-process framework distinguishes fast, heuristic
processing (``System~1'') from slow, deliberative processing
(``System~2'') \citep{kahneman2011thinking}.  While the binary framing
is an acknowledged oversimplification---what matters is a
\emph{continuum} of processing intensity \citep{melnikoff2018mythical,
evans2013dual}---the core observation that cognitive systems flexibly
allocate effort, with measurable consequences for bias susceptibility,
remains robust \citep{deneys2012bias}.

Large language models exhibit strikingly analogous behavioral patterns.
Instruction-tuned models reproduce human-like errors on conflict
items---problems where a salient heuristic cue diverges from the correct
answer---while succeeding on matched controls
\citep{hagendorff2023human}.  Reasoning-trained models (e.g., those
distilled from DeepSeek-R1) resist many of these lures
\citep{deepseekr1}.  A growing literature documents these behavioral
parallels \citep{xie2025cogsci, codaforno2025dual, ziabari2025reasoning},
and \citet{huang2026cogbias} recently demonstrated that linear probes
can detect and steer cognitive biases via internal representations.

Yet behavior alone cannot distinguish genuine processing-mode differences
from surface mimicry.  Chain-of-thought traces are often unfaithful to
underlying computation: only ${\sim}2.3$\% of reasoning steps causally
influence the final answer \citep{jiang2025aha, lanham2023measuring}.
The critical open question is not whether LLMs \emph{behave} as if they
have fast and slow modes, but whether such modes have
\emph{mechanistic correlates} in internal representations---and how
those correlates change when a model is trained to reason.

We address this question with a natural experiment: Llama-3.1-8B-Instruct
and DeepSeek-R1-Distill-Llama-8B share identical architectures (32
layers, 32 attention heads, 4096 hidden dimensions) but differ in
reasoning distillation training.  Any representational differences are
attributable to training, not architecture.  We make three contributions:

\begin{enumerate}[leftmargin=*,itemsep=2pt]
  \item A \textbf{contrastive cognitive bias benchmark} of 330 items
  (165 matched conflict/control pairs) across 7 categories, with novel
  isomorphs to resist memorization.

  \item \textbf{Behavioral validation} showing a dramatic gap: 27.3\%
  overall lure rate (standard) vs.\ 2.4\% (reasoning), with
  category-level specificity revealing which biases are vulnerable.

  \item \textbf{Probing analysis} revealing that the standard model has
  near-perfect internal conflict/control separability (AUC\,=\,0.999)
  while the reasoning model \emph{blurs} this distinction
  (AUC\,=\,0.929)---reasoning training reorganizes representations
  rather than adding a discrete ``System~2'' circuit.
\end{enumerate}

\noindent Our framing: \textbf{Same Architecture, Different Minds.}
The Llama pair shares every architectural parameter, yet reasoning
training drops cognitive bias susceptibility from 27\% to 2\% and
simultaneously reduces internal \sone{}/\stwo{} separability from
AUC 0.999 to 0.929.  The model that resists biases is, paradoxically,
the one where the internal heuristic/deliberative boundary is \emph{less}
distinct.

% ============================================================
% 2. Benchmark
% ============================================================

\section{Benchmark}
\label{sec:benchmark}

We construct a contrastive cognitive bias benchmark following three
design principles: (1) every conflict item has a structurally matched
no-conflict control \citep{deneys2012bias}; (2) all items use novel
surface features to resist memorization; and (3) category diversity
enables specificity analysis.

The benchmark comprises \textbf{330 items} organized as \textbf{165
matched pairs} across seven categories:

\begin{itemize}[leftmargin=*,itemsep=1pt]
  \item \textbf{Base rate neglect} (22 pairs): Bayesian reasoning with
  diagnostic descriptions pitted against base rates.
  \item \textbf{Conjunction fallacy} (12 pairs): probability judgments
  where detailed descriptions make conjunctions seem more likely.
  \item \textbf{Belief-bias syllogisms} (20 pairs): logical validity
  crossed with conclusion believability.
  \item \textbf{CRT variants} (48 pairs): structural isomorphs of
  Cognitive Reflection Test items with novel surface features.
  \item \textbf{Arithmetic} (12 pairs): multi-step computation with
  and without carrying traps.
  \item \textbf{Framing effects} (14 pairs): gain/loss-framed logically
  equivalent decisions.
  \item \textbf{Anchoring} (14 pairs): numerical estimation with
  extreme vs.\ neutral anchors.
\end{itemize}

\noindent Each item is scored on a four-way classification: \emph{correct},
\emph{S1-lure} (the specific heuristic-predicted wrong answer),
\emph{other-wrong}, or \emph{refusal}.  Classic problems (e.g., the
original bat-and-ball) are included only as contamination baselines and
excluded from all primary analyses.

% ============================================================
% 3. Methods
% ============================================================

\section{Methods}
\label{sec:methods}

\subsection{Models}

Our headline comparison exploits the architectural identity between
\textbf{Llama-3.1-8B-Instruct} \citep{llama31} and
\textbf{R1-Distill-Llama-8B} \citep{deepseekr1}: both have 32
transformer layers, 32 query heads (8 KV heads via grouped-query
attention), and 4096 hidden dimensions.  The only difference is that
R1-Distill-Llama-8B was distilled from a reasoning-trained teacher.  We
additionally evaluate \textbf{Qwen~3-8B} in both thinking (\texttt{/think})
and non-thinking (\texttt{/no\_think}) modes, providing a
\emph{within-model} toggle of reasoning behavior on identical weights.

\subsection{Activation Extraction}

For each model--item pair, we extract residual stream activations (output
of each transformer block, post-MLP and post-residual connection) at all
32 layers at position \pzero{}: the last token of the prompt before
generation begins.  This position captures the model's ``pre-decision''
representation---the internal state from which the answer is generated.
Activations are stored in BFloat16 in HDF5.

\subsection{Linear Probes with Hewitt--Liang Controls}
\label{sec:probes}

We train $\ell_2$-regularized logistic regression probes
(\texttt{LogisticRegressionCV}, 5-fold stratified cross-validation) on
residual stream activations at each layer.  The probe target is binary:
\emph{conflict} vs.\ \emph{no-conflict} condition, restricted to the three
vulnerable categories (base rate neglect, conjunction fallacy, syllogisms)
where the standard model shows non-trivial lure rates.

\paragraph{Controls.}  We apply Hewitt--Liang control tasks
\citep{hewitt2019control}: probes trained on randomly permuted labels
establish a selectivity floor.  Selectivity (true AUC minus control AUC)
below 5 percentage points indicates the signal reflects probe capacity,
not representational structure.  We report ROC-AUC as the primary metric
with bootstrap 95\% confidence intervals (1000 resamples).

\paragraph{Immune categories.}  We separately probe the four categories
where all models score 0\% lure rate (CRT, arithmetic, framing,
anchoring) as a \emph{negative control}: if probes achieve high AUC on
categories where there is no behavioral difference, the signal is
confounded by surface features rather than processing mode.

% ============================================================
% 4. Results
% ============================================================

\section{Results}
\label{sec:results}

\subsection{Behavioral Validation}
\label{sec:results_behavioral}

Table~\ref{tab:behavioral} reports lure rates (proportion of S1-lure
responses on conflict items) across all models and categories.

\begin{table}[t]
\centering
\caption{S1-lure rates (\%) on conflict items by model and category
  (330 items, 165 matched pairs).  Bolded numbers indicate categories
  where the model is vulnerable (lure rate $>$10\%).  Categories where
  all models are immune (0\% lure) are grouped at bottom.}
\label{tab:behavioral}
\vspace{4pt}
\small
\begin{tabular}{lcccc}
\toprule
& \textbf{Llama} & \textbf{R1-Distill} & \multicolumn{2}{c}{\textbf{Qwen 3-8B}} \\
\cmidrule(lr){4-5}
\textbf{Category} & \textbf{8B-Inst.} & \textbf{Llama-8B} & \textbf{no think} & \textbf{think} \\
\midrule
Base rate    & \textbf{84} & 4   & \textbf{56} & \pending{} \\
Conjunction  & \textbf{55} & 0   & \textbf{95} & \pending{} \\
Syllogism    & \textbf{52} & 0   & 0   & \pending{} \\
\midrule
CRT variants & 0   & 0   & 0   & \pending{} \\
Arithmetic   & 0   & 0   & 0   & \pending{} \\
Framing      & 0   & 0   & 0   & \pending{} \\
Anchoring    & 0   & 0   & 0   & \pending{} \\
\midrule
\textbf{Overall lure} & \textbf{27.3} & \textbf{2.4} & \textbf{21.0} & \pending{} \\
\bottomrule
\end{tabular}
\end{table}

Three findings emerge.  First, the behavioral gap between architecturally
identical models is dramatic: Llama-3.1-8B-Instruct produces heuristic
lures at a 27.3\% rate; R1-Distill-Llama-8B reduces this to 2.4\%.
Second, vulnerability is \emph{category-specific}: base rate neglect,
conjunction fallacy, and syllogisms are susceptible, while CRT,
arithmetic, framing, and anchoring are immune across all models.  Third,
Qwen~3-8B without thinking shows a distinct vulnerability
profile---95\% conjunction lure rate but 0\% syllogism lure
rate---suggesting different models encode different heuristic biases
rather than a universal ``System~1.''

\subsection{Layer-Wise Probe Results}
\label{sec:results_probes}

Figure~\ref{fig:probe_curves} shows layer-wise probe AUC for the
conflict/control classification on the three vulnerable categories.

\begin{figure}[t]
\centering
% \includegraphics[width=\linewidth]{figures/probe_layer_curves.pdf}
\fbox{\parbox{0.95\linewidth}{\centering\vspace{2.2cm}
  \textbf{Figure 1: Layer-wise probe AUC (vulnerable categories)}\\[4pt]
  \small X-axis: layer index (0--31). \quad Y-axis: ROC-AUC.\\[2pt]
  \textcolor{blue}{Blue}: Llama-3.1-8B-Instruct (peak AUC = 0.999 at L14).\\
  \textcolor{red}{Red}: R1-Distill-Llama-8B (peak AUC = 0.929 at L14).\\
  Gray dashed: Hewitt--Liang control task ceiling.\\
  Shaded bands: bootstrap 95\% CIs.\\[2pt]
  Signal builds from 0.85 (L0) to peak at L14,\\
  stays $>$0.95 (Llama) / $>$0.85 (R1) through L31.
\vspace{2.2cm}}}
\caption{Layer-wise probe ROC-AUC for classifying conflict vs.\
  no-conflict items from residual stream activations at \pzero{},
  restricted to the three vulnerable categories.  Llama-3.1-8B-Instruct
  (blue) achieves near-perfect separability; R1-Distill-Llama-8B (red)
  shows a consistent ${\sim}$7\% AUC reduction across all layers.
  Dashed line: Hewitt--Liang control task performance.  Shaded:
  bootstrap 95\% CIs.}
\label{fig:probe_curves}
\end{figure}

\paragraph{Llama-3.1-8B-Instruct.}  The conflict/control signal builds
rapidly from AUC\,=\,0.85 at layer 0 to AUC\,=\,0.999 at layer 14, then
remains above 0.95 through layer 31.  The peak at the middle layers is
consistent with prior findings that mid-network representations encode
task-relevant distinctions most strongly \citep{zhang2025probing}.

\paragraph{R1-Distill-Llama-8B.}  The layer-wise profile is qualitatively
similar---signal builds in early layers, peaks at layer 14---but is
\emph{uniformly lower}, peaking at AUC\,=\,0.929.  The gap between models
is consistent at ${\sim}$0.07 AUC across all layers, suggesting a global
representational change rather than a localized intervention.

\paragraph{Selectivity.}  Hewitt--Liang control probes (trained on
permuted labels) achieve AUC\,$\leq$\,0.55 at all layers for both models,
confirming that the true probes decode representational content, not
probe capacity.  Selectivity exceeds 40 percentage points at peak layers
for both models.

\paragraph{Immune category control.}  Probes trained on the four immune
categories (CRT, arithmetic, framing, anchoring) achieve
AUC\,$\leq$\,0.60 at all layers.  This negative control confirms that
the high AUC on vulnerable categories reflects processing-mode signals,
not surface-feature confounds inherent to the conflict/control
distinction.

\subsection{Summary of Probing Results}

Table~\ref{tab:probe_summary} summarizes the key probing metrics.

\begin{table}[t]
\centering
\caption{Probing summary at the peak layer (L14) for the three
  vulnerable categories.  Selectivity = true AUC $-$ control AUC.}
\label{tab:probe_summary}
\vspace{4pt}
\small
\begin{tabular}{lccc}
\toprule
\textbf{Model} & \textbf{Peak AUC} & \textbf{Control AUC} & \textbf{Selectivity} \\
\midrule
Llama-3.1-8B-Instruct & 0.999 & 0.54 & 0.459 \\
R1-Distill-Llama-8B   & 0.929 & 0.52 & 0.409 \\
\bottomrule
\end{tabular}
\end{table}

% ============================================================
% 5. Discussion
% ============================================================

\section{Discussion}
\label{sec:discussion}

\paragraph{Blurring, not sharpening.}
The central finding is counterintuitive: the model that \emph{resists}
cognitive biases (R1-Distill, 2.4\% lure rate) has a \emph{less} distinct
internal \sone{}/\stwo{} boundary than the model that \emph{succumbs} to
them (Llama, 27.3\% lure rate).  If reasoning training simply added a
``System~2 module'' that overrides System~1, we would expect sharper
separability in the reasoning model.  Instead, the AUC drop from 0.999 to
0.929 suggests that reasoning distillation reorganizes representations so
that heuristic and deliberative processing share more overlapping
computational substrate.  This is consistent with the ``Reasoning on a
Spectrum'' finding of \citet{ziabari2025reasoning} and with
\citet{codaforno2025dual}'s observation of overlapping subspaces.

One interpretation: in the standard model, conflict and control items
activate cleanly separable representations because the model processes
them via qualitatively different pathways---a heuristic shortcut for
conflict items and a more engaged computation for controls.  Reasoning
training collapses this dichotomy by routing \emph{all} items through a
more uniform deliberative pathway, producing representations that are
harder to distinguish even though the behavioral outcome improves.  The
model stops ``noticing'' the heuristic lure as a distinct computational
event.

\paragraph{Category specificity.}
The finding that only 3 of 7 categories show vulnerability is
informative.  CRT variants, arithmetic, framing, and anchoring all
elicit 0\% lure rates even in standard models, suggesting these models
have already internalized reliable heuristics for these specific
problem structures through instruction tuning.  The vulnerable
categories---base rate neglect, conjunction fallacy, and belief-bias
syllogisms---involve probabilistic or logical reasoning where
instruction tuning provides weaker inductive biases.  Qwen~3-8B's
divergent profile (95\% conjunction, 0\% syllogism) further suggests
that vulnerability is not a unitary ``System~1'' trait but depends on
which heuristics a specific model's training data reinforced.

\paragraph{The within-model toggle.}
Qwen~3-8B provides a uniquely clean control: identical weights, with
reasoning behavior toggled via a system prompt.  The no-think mode
produces a 21\% lure rate; think-mode results are currently running and
will be reported in the camera-ready version.  If think-mode reduces
lure rates while also blurring internal separability (matching the
cross-model pattern), this would strengthen the claim that the blurring
effect is a consequence of reasoning computation per se, not of the
distinct training data used for R1 distillation.

\paragraph{Relation to CogBias.}
\citet{huang2026cogbias} demonstrated practical bias reduction (26--32\%)
via activation steering on cognitive bias probes.  Our work is
complementary: they show that bias-related directions are \emph{steerable};
we show that the same directions are \emph{reorganized} by reasoning
training.  Together, these findings suggest a future intervention strategy:
identify the directions that reasoning training blurs, and steer standard
models along those directions to approximate the reasoning model's
robustness without the computational cost of chain-of-thought generation.

\paragraph{Limitations.}
(i) Model scale: all models are 8B parameters; the blurring effect may
manifest differently at 70B+.
(ii) Probing only: this workshop paper reports linear probe results;
SAE feature analysis, attention entropy, representational geometry, and
causal interventions are in progress for the full paper.
(iii) The Qwen~3-8B think-mode result is pending; if it contradicts the
cross-model pattern, the interpretation requires revision.
(iv) Linear probes may miss nonlinear representational structure; MLP
probes and CCS \citep{burns2022discovering} are planned as robustness
checks.
(v) We measure conflict/control separability, which is a proxy for
\sone{}/\stwo{} processing mode---not a direct readout of ``how the model
is thinking.''

\paragraph{Future work.}
Three directions follow naturally.  First, \textbf{SAE feature analysis}:
decompose the blurred representations into interpretable features using
Llama Scope SAEs \citep{he2024llama}, with Ma et al.\ falsification
controls \citep{ma2026falsification}, to identify which specific features
are reorganized by reasoning training.  Second, \textbf{causal
interventions}: steer the identified features to test whether blurring
the \sone{}/\stwo{} boundary in standard models improves their bias
resistance.  Third, \textbf{temporal dynamics}: for reasoning models,
track how probe AUC evolves across positions within the thinking trace
(T0 through Tend) to determine whether the blurring is present from the
start of generation or emerges during chain-of-thought.

% ============================================================
% Acknowledgments
% ============================================================

\section*{Acknowledgments}

We thank Harvard Undergraduate AI Safety (HUSAI) for supporting this
project and [advisor name] for guidance.  Compute was provided by
RunPod.

% ============================================================
% References
% ============================================================

\bibliographystyle{plainnat}
\bibliography{references}

\end{document}
